Une compagnie aérienne mène une enquête de satisfaction auprès de ses clients. Les deux questions posées lors de cette enquête sont les suivantes :\\
 \og Avez-vous acheté votre billet d'avion en agence de voyage ou par internet ? \fg{}\\
	\og Êtes-vous satisfait des services proposés par notre compagnie ? \fg{}.

	\bigskip 
	Les résultats sont donnés dans le tableau suivant :
	
	\begin{center}
		\begin{tblr}{hlines,vlines,colspec={X[1.4,l]X[c]X[c]X[c]},cells={m},row{1}={font=\bfseries},column{1}={font=\bfseries}}
			& Billet acheté par internet & Billet acheté en agence & Total \\
			Satisfait & 720 & 180 & 900 \\
			Pas satisfait & 80 & 20 & 100 \\
			Total & 800 & 200 & 1\,000 \\
		\end{tblr}
	\end{center}
	
	On choisit au hasard un client ayant répondu aux deux questions. On considère les deux événements suivants :
	\begin{itemize}[label={}]
		\item $S$ : \og Le client est satisfait \fg{} ;
		\item $I$ : \og Le client a acheté son billet par internet \fg{}.
	\end{itemize}
	
	\begin{enumerate}
		\item Calculer la probabilité qu'un client soit satisfait.
		\item Calculer la probabilité de l'événement $S\cap I$. \\
		Interpréter le résultat à l'aide d'une phrase.
		\item Calculer $\text{P}_{S}(I)$.
		\item Les deux événements $I$ et $S$ sont-ils indépendants ? Justifier la réponse.
		\item Le service marketing de la compagnie aérienne affirme que $90~\%$ des clients qui ont acheté leur billet en agence sont satisfaits. Est-ce le cas ? Justifier la réponse.
	\end{enumerate}
